\begin{table}[htbp]
\centering
\caption{Sample construction}
\label{tab:waterfall}
\small
\begin{threeparttable}
\begin{tabular}{@{}lrr@{}}
\toprule
Step & Stock-days & \% of attempted \\
\midrule
Stock-days attempted & 2,501 & 100.0 \\
Survived the pre-screen & 2,493 & 99.7 \\
First trade by 9:10 & 2,488 & 99.5 \\
Opening price above \textyen 200 & 2,493 & 99.7 \\
More than 20 continuous-session trades & 2,490 & 99.6 \\
Tick a power of ten all day & 2,234 & 89.3 \\
Passing all four filters & 2,229 & 89.1 \\
Stock qualifies on more than half the sample & 2,182 & 87.2 \\
\bottomrule
\end{tabular}
\begin{tablenotes}[flushleft]\footnotesize
\item The four filters are Ohta's, applied at tick level. The pre-screen is a cheap four-column read that rejects stock-days which cannot possibly qualify -- overwhelmingly those with too few trades -- and changes nothing about the final sample. The filters are not nested, so their individual pass counts do not multiply to the joint count.
\end{tablenotes}
\end{threeparttable}
\end{table}
